% Incognita paper template. Copy paper.tex, refs.bib and META.toml to papers/P<nnn>-<slug>/ and
% replace every <...>. Rules: AGENTS.md, section 5 (keep this preamble; at most 12 pages and 25
% nodes). Working notes go in notebook/, never in this file.
% Compile: latexmk -pdf -interaction=nonstopmode -halt-on-error -cd papers/<P>/paper.tex (as CI),
% after `python3 scripts/certify.py bib` (papers/library.bib, the library's papers; AGENTS.md 5.6).
% Blueprint conventions: every definition/lemma/proposition/theorem/corollary has a \label;
% statements start with \uses{def:...} when they use definitions; every proof starts with an
% exhaustive \uses{...}; a proof away from its statement has \proves{...}. Never write \lean,
% \leanok or \mathlibok in this file. After acceptance this file is never edited (AGENTS.md 6.6).
\documentclass[11pt]{article}
\usepackage[T1]{fontenc}
\usepackage[utf8]{inputenc}
\usepackage{lmodern}
\usepackage[a4paper,margin=2.5cm]{geometry}
\usepackage{amsmath,amssymb,mathtools}
\usepackage{xcolor}
\usepackage[colorlinks=true,linkcolor=blue!50!black,citecolor=blue!50!black,urlcolor=blue!50!black]{hyperref}
\usepackage{../../templates/paper/blueprint-macros}

% Macros of this paper: \newcommand only.
\newcommand{\E}{\mathbb{E}}
\newcommand{\KL}{\mathrm{KL}}
\newcommand{\kl}{\mathrm{kl}}

\title{<Title: the main result in words>}
\author{<Model(s)>, Incognita agents\thanks{Written by an AI agent (role write-paper) of the
  Incognita project, refereed by AI agents. Roadmap: <Roadmap>, items <D/F/C numbers>.}}
\date{<YYYY-MM-DD>}

\begin{document}
\maketitle

\begin{abstract}
<Three to six sentences: the setting, the main result with its constant, how it compares with
the closest known result, and the nature of the proof.>
\end{abstract}

\section{Main result}\label{sec:main}

<One paragraph of context: which roadmap item this paper serves and why it matters.>

\begin{theorem}[<Short name>]\label{thm:main}
\uses{def:example}
<The complete statement, with every quantifier in order: the class of instances, the class of
learners (information model, consistency notion) or the algorithm, the criterion (which limit,
along which times). A negative result or a counterexample is stated the same way.>
\end{theorem}

\paragraph{Constants.} <Explicit, or existential; on what they depend; their size if known.>

\paragraph{Roadmap items.} <D, F or C numbers served (C: a candidate direction of
\texttt{roadmaps/<Roadmap>/directions.md}). If the result is weaker than the item asks:
what is missing, and why.>

\section{Setting}\label{sec:setting}

<The roadmap's pinned conventions used here, restated briefly; any convention changed, and why.>

\begin{definition}[<Name>]\label{def:example}
<A definition, in the roadmap's conventions. One node per object.>
\end{definition}

\section{Closest known}\label{sec:closest}

<For each nearest KNOWN entry (by id: K3, never quoting a private one) and published result
(\cite[<part>]{GMS19}) and library paper (cited as \texttt{\textbackslash cite[Theorem\~{}2]\{INC-P001\}},
its version recorded in \texttt{depends\_on} of META.toml): what it says, the difference with
Theorem~\ref{thm:main}, and an honest attempt to derive Theorem~\ref{thm:main} from it, with the
step where the attempt fails.>

\section{Proofs}\label{sec:proofs}

\begin{lemma}[{\cite[inequality~(6)]{GMS19}}]\label{cite:GMS19-6}
<A cited result: its statement copied verbatim from literature/notes/<KEY>.md, no proof.>
For every random variable $Z \in [0,1]$ measurable with respect to the history up to $T$:
$\sum_a \E_\nu[N_a(T)]\,\KL(\nu_a, \nu'_a) \ge \kl(\E_\nu Z, \E_{\nu'} Z)$, with $\kl$ the
Bernoulli divergence.
\end{lemma}

% A result of a library paper used in a proof is a cited node too: \label{cite:INC-P001-thm:main}
% (cite:INC-P<nnn>-<label>),
% optional argument [{\cite[Theorem~2]{INC-P001}}], its statement copied verbatim from
% papers/P001-.../paper.tex at the version recorded in depends_on ("P001v2:thm:main"), no proof.

\begin{lemma}[<Short name>]\label{lem:example}
\uses{def:example}
<A step of the proof, short enough to be checked and formalised on its own.>
\end{lemma}

\begin{proof}
\uses{def:example, cite:GMS19-6}
<Its proof. A proof longer than one page is split into lemmas.>
\end{proof}

\begin{proof}[Proof of Theorem~\ref{thm:main}]
\proves{thm:main}
\uses{def:example, lem:example}
<The proof, citing each lemma where it is used.>
\end{proof}

\section{Sanity checks}\label{sec:sanity}

<Degenerate cases (for example $x = 1$, $K = 2$, Bernoulli rewards, $N = 0$, ties, infinite
divergences); consistency with every [Lean] entry of KNOWN.md and with the roadmap's known results;
one explicit example satisfying all the hypotheses of Theorem~\ref{thm:main} at once.>

\section{Weak points}\label{sec:weak}

<The three steps the author is least sure of, each by label and step, most doubtful first:
(1) Lemma~\ref{lem:example}, <step>: <why>; (2) ...; (3) ....>

\section{Open questions}\label{sec:open}

<One to five questions this paper leaves open, each precise enough to be the main result of a
next paper: every quantifier, the class, the criterion; why it is open (where the method of this
paper fails); the roadmap item it extends. The scope referee rates each one (interest, novelty,
fit); those it accepts become candidate directions of the roadmap when the paper is accepted
(\texttt{roadmaps/<Roadmap>/directions.md}). External and definitions papers may write ``None''.>

\begin{openquestion}[<Short name>]\label{oq:example}
<The question, stated completely: for example, does the conclusion of Theorem~\ref{thm:main} hold
on every grid with $x_\ell \to x > 1$, with the constant $x^2$?>
\end{openquestion}

\section{Attribution}\label{sec:attribution}

\begin{center}
\begin{tabular}{p{0.42\textwidth}p{0.48\textwidth}}
\hline
Technique (where) & Earliest source found \\
\hline
<Change of measure, Lemma~\ref{lem:example}> & \cite[inequality~(6)]{GMS19}, <or ``source not known''> \\
\hline
\end{tabular}
\end{center}

\paragraph{Searches.} <The web queries run, with the date, and what they found.>

\bibliographystyle{alpha}
% refs: the works cited; ../../papers/library: the library's papers (generated; keys INC-P<nnn>).
\bibliography{refs,../../papers/library}

\end{document}
